\documentclass[11pt]{letter}
\usepackage[margin=1in]{geometry}
\usepackage{times}

\signature{Jeff Gore (Corresponding Author) \\ Department of Physics \\ Massachusetts Institute of Technology \\ Cambridge, MA, USA \\\\ Jinyeop Song \\ Jiliang Hu}

\address{Department of Physics \\ Massachusetts Institute of Technology \\ Cambridge, MA 02139, USA}

\begin{document}

\begin{letter}{Editorial Office \\ [Journal Name]}

\opening{Dear Editor,}

We are pleased to submit our manuscript, ``Interspecies Interactions Drive Community-Level Selection in Microbial Coalescence,'' for consideration in [Journal Name].

A fundamental and unresolved question in ecology is whether communities function as cohesive, integrated units of selection---the ``superorganism'' view championed by Clements---or merely as loose collections of independent species, as proposed by Gleason's individualistic hypothesis. Community coalescence, the mixing of previously isolated communities, provides a natural testing ground for these competing paradigms: cohesive communities should persist as units during coalescence, while loose assemblages should sort independently. Yet empirical studies have yielded conflicting results. Some systems exhibit clear community-level selection, while others show independent species sorting, with no theoretical framework to reconcile this discrepancy.

In this work, we combine generalized Lotka-Volterra modeling with high-throughput bacterial microcosm experiments involving 54 isolates spanning 29 families and three phyla to demonstrate that \textit{interspecies interaction strength} is the governing parameter that determines the level of selection during coalescence. Our simulations predict, and our experiments confirm, that weak interactions lead to independent species sorting (Gleasonian dynamics), whereas moderate-to-strong interactions drive correlated species persistence and community-level selection (Clementsian dynamics). We validate this experimentally by manipulating nutrient concentrations to tune interaction strength across three regimes, toggling the system between species-level and community-level selection in quantitative agreement with our theoretical predictions.

Beyond identifying interaction strength as the control parameter, we uncover two mechanistically distinct regimes of community-level selection: an \textit{emergent} regime in which collective multi-species dynamics shape coalescence outcomes, and a \textit{top-down} regime in which dominant species determine the winning community through environmental modification (specifically, pH). We further demonstrate that these interaction-dependent patterns generalize from synthetic consortia to complex communities derived from six diverse natural environments.

CCrucially, our framework reconciles conflicting observations in the recent literature by identifying interaction strength as the "tuning knob" for selection. This provides a unified physical perspective on community assembly that explains why different environments favor different selection regimes.

Given the broad implications for theoretical ecology, microbiome engineering, and our understanding of biological complexity, we believe this manuscript will be of significant interest to the readership of [Journal Name]..

This manuscript has not been published and is not under consideration for publication elsewhere. We have no conflicts of interest to disclose.

\closing{Sincerely,}

\end{letter}
\end{document}
